% <<<PDFPAGE 24>>>
\chapter{KV 缓存优化}

\section{KV 缓存的显存压力}

如第 3 章所示，单个长上下文请求的 KV 缓存就可能占用几十 GB 显存。服务系统如果同时处理数百个请求，直接按最大序列长度预分配 KV 缓存会浪费大部分已分配的显存，因为请求的实际长度难以预知。这也严重限制了批大小。2023 年以前，这一直是大语言模型服务中一个尚未解决的主要问题。

\section{PagedAttention 与 KV 缓存的虚拟内存}

PagedAttention（Kwon 等，2023）是 vLLM 的主要创新。它把操作系统的虚拟内存（virtual memory）概念用于 KV 缓存管理。

PagedAttention 将 KV 缓存划分成大小固定的\textbf{块}，也称为页，每块容纳 \(B\) 个词元。它用一张\textbf{块表}（block table），把每个请求在 KV 缓存中的逻辑位置映射到显存中的物理块。这些物理块可以不连续，因而无须给每个请求划出一整段连续的显存。

\begin{conceptbox}{PagedAttention 的主要特性}
\begin{itemize}
  \item \textbf{无须预分配}。随生成过程按需分配块。
  \item \textbf{无内部碎片}。每个请求浪费的显存至多相当于 \(B-1\) 个词元。
  \item \textbf{写时复制}（copy-on-write）。使用相同前缀（例如系统提示）的多个请求可共享物理块，直至各自的生成路径分岔。
  \item \textbf{可灵活淘汰}。可将一个请求的块换出至 CPU 内存或磁盘，为另一个请求腾出空间。
\end{itemize}
\end{conceptbox}

\translatornote{“无内部碎片”为原文用语。按本条所述，一个请求的末块若未填满，仍可能有至多 \(B-1\) 个词元槽位空置。}

PagedAttention 需要修改注意力内核，使它在计算注意力时收集不连续的块。实现起来并不简单，vLLM 为此采用了定制 CUDA 内核。

\section{前缀缓存（提示缓存）}

在许多生产工作负载中，数千个请求会共享一段较长的前缀，例如系统提示、文档或少样本示例（few-shot example）。若没有缓存，每个请求都需要从头为这段共享前缀计算 KV 缓存。

\textbf{前缀缓存}也称为提示缓存（prompt caching）。它保存常用前缀的 KV 缓存，供不同请求复用。vLLM 和 SGLang 实现了这一机制，Anthropic 和 OpenAI 也将它作为功能提供。当前缀相对于用户提示较长时，预填充成本可降低 90\% 以上。

% <<<PDFPAGE 25>>>
\section{StreamingLLM 与注意力汇聚现象}

流式对话和长时间运行的智能体等应用要求上下文可以无限延长。若保留所有词元的 KV 缓存，缓存就会无界增长。朴素的滑动窗口注意力（sliding window attention）只保留最近 \(W\) 个词元的 KV 缓存，结果会使模型表现急剧下降。原因在于大语言模型会形成\textbf{注意力汇聚点}（attention sinks）。无论语义内容是什么，序列最前面的几个词元都会在各层获得异常高的注意力分数。删去它们会破坏模型的内部表示。

\textbf{StreamingLLM}（Xiao 等，2023）保留少数起始词元作为汇聚点，通常保留 4 个，同时保留最近词元的滑动窗口。这样就能以有界的显存支持任意长度的大语言模型推理。
\[
  \text{保留 KV 缓存的词元}
    =\text{汇聚点词元}\cup\text{最近窗口内的词元}
\]

\section{KV 缓存量化}

KV 缓存本身可以独立于模型权重进行量化。每一步注意力计算只读取一次缓存，量化误差只影响当前解码步骤，不会逐步累积。在 TensorRT-LLM、vLLM 等框架中，INT8 KV 缓存量化已成为常规做法，可使 KV 缓存的显存占用降至约一半。更激进的方案会对 KV 缓存采用 INT4，甚至 FP8，并谨慎地按注意力头或通道选取缩放尺度。

\begin{interview}
\question{1}{请解释 PagedAttention。它与 KV 缓存的连续预分配方式有何区别？}
传统的 KV 缓存管理为每个请求预分配一段连续显存，大小按可能的最大序列长度确定。大多数请求远短于这个最大长度，因此显存浪费严重。此外，释放一块显存后，留下的空隙可能装不下新的块，从而造成外部碎片（external fragmentation）。PagedAttention 将 KV 缓存的显存划分成固定大小的块，每块约容纳 16 个词元，并为每个请求维护一张从逻辑位置到物理块的映射表。生成新词元时按需分配块，请求完成时立即释放块。这样省去了预分配造成的浪费，将每个请求的碎片限制在至多 \(B-1\) 个词元的空间内，还能以写时复制方式跨请求共享公共前缀块。在相同显存容量下，批大小通常可达到原来的 \(2\) 至 \(4\) 倍。

\question{2}{前缀缓存如何工作？哪类应用从中受益最多？}
前缀缓存保存一段词元前缀（例如系统提示）已算出的 KV 缓存，供拥有同一前缀的后续请求复用。系统用哈希表（hash map）把前缀词元序列或内容的哈希值映射到 KV 缓存的物理块。新请求到达时，服务系统检查开头的词元是否与某个已缓存前缀匹配。若匹配，就可以跳过这些词元的预填充。受益最多的应用包括三类。（1）\textbf{使用固定系统提示的 API 产品}。例如客服机器人使用长达 2000 个词元的系统提示，缓存后每条用户消息都能复用此前缀。（2）\textbf{多轮对话}。后续各轮可以复用此前各轮的 KV 缓存。（3）\textbf{批量文档处理}。
% <<<PDFPAGE 26>>>
例如，有 1000 个问题指向同一份文档，所有请求可以共享该文档的 KV 缓存。在实际应用中，预填充阶段的时延可降低 50\% 至 90\%。

\question{3}{什么是注意力汇聚现象？为什么朴素的滑动窗口注意力会在 StreamingLLM 所处理的流式推理中失效？}
当大语言模型用固定长度的序列训练时，会学到把大量注意力分配给每个序列开头的几个词元。这些词元称为注意力汇聚点。为了使 softmax 归一化，必须分配的多余注意力概率质量会集中在这些起始词元上。这些词元的语义并不重要，只是承担结构上的锚定作用。朴素的滑动窗口注意力会丢弃最早的词元，以维持固定大小的 KV 缓存。一旦窗口滑过开头，汇聚点也会被移除。失去汇聚点后，注意力的归一化被扰乱，剩余词元的 softmax 分数与训练时的分布差异很大，导致输出质量急剧下降。StreamingLLM 始终在 KV 缓存中保留最初 4 个汇聚点词元和最近窗口中的词元，从而恢复正常的注意力行为。

\question{4}{如果系统所处理的请求长度差异很大，你会如何设计 KV 缓存淘汰策略？}
我会综合以下信息，为 KV 缓存淘汰（cache eviction）的候选块确定优先级。（1）\textbf{访问时间}。按最近最少使用（Least Recently Used，LRU）原则优先淘汰最久未被访问的块。（2）\textbf{序列结束的概率}。结合请求当前的生成位置和典型长度分布，估算它离结束还有多远，优先淘汰预计很快会结束的请求所占用的块。（3）\textbf{块的可共享性}。保护由多个请求共享的块，例如前缀缓存块，因为复用可为多个请求节省开销。（4）\textbf{淘汰成本}。优先淘汰提示较短、KV 缓存重算成本低的请求的块，尽量保留预填充耗时长的请求的块。CPU 内存较充裕，我还会把 CPU DRAM 作为第二级存储，在直接淘汰前先把块换出到那里。vLLM 的抢占（preemption）机制就是一个实际例子。当显存紧张时，它暂停请求并重新调度。

\question{5}{请解释 KV 缓存量化。KV 缓存能否比模型权重采用更激进的量化？为什么？}
KV 缓存量化以 INT8、INT4 或 FP8 等低于 FP16 精度的格式存储键和值，也就是注意力缓存。支持比权重采用\textbf{更激进}的量化有几个理由。（1）每个注意力步骤只读取一次 KV 缓存的值，各步之间不会累积误差，因为每次解码都重新从缓存读取。（2）注意力本身含有归一化操作（softmax），因而对键的量化误差有一定的耐受能力。（3）相较于权重矩阵，KV 缓存在不同头与不同层之间有更多冗余。另一方面，KV 缓存来源于激活值，而激活值具有离群现象。和激活量化中遇到的问题一样，这使缓存比权重更难量化。实践中，INT8 KV 缓存的做法已经成熟，质量损失很小。INT4 KV 缓存可以采用逐头或逐词元的缩放，但需要仔细实现。
\end{interview}
